\chapter{消化与营养物质}\label{ch:g7:digestion-and-nutrients}

\Cref{ch:g4:journey-of-food}中，一口面包顺着\omterm{def:g4:journey-of-food:digestion}{消化}管向下
走，并许诺“\omterm{def:g4:journey-of-food:digestion}{消化液}会把它拆解开”。
本章我们推开“加工车间”的大门：\omterm{def:g4:journey-of-food:digestion}{消化液}究竟做了什么，
它们产生了哪些\omterm{def:g7:digestion-and-nutrients:families}{营养物质}，以及\omterm{def:g4:journey-of-food:tract}{小肠}壁如何
完成人体内规模最大的一次
跨界穿越？一口面包又回来了，这次
借助更精密的仪器再看一看。

\section{消化液的工作}

\begin{proposition}[消化是化学拆解]\label{prop:g7:digestion-and-nutrients:chemical}
咀嚼和搅拌只是把食物\emph{切小}；真正的
拆解是化学过程。\emph{消化
液}\index{消化液}——\omterm{def:g4:journey-of-food:tract}{唾液}、胃液、
注入\omterm{def:g4:journey-of-food:tract}{小肠}的那些\omterm{def:g4:journey-of-food:digestion}{消化液}——
都含有起作用的物质，能把食物的大分子
切成小分子：淀粉变成糖，肉和蛋的
构建物质变成其小单位，脂肪也分解成各自的小单位。
每种\omterm{def:g4:journey-of-food:digestion}{消化液}都有自己的专长，各自
在所处位置的\omterm{def:g6:exploring-our-environment:conditions}{条件}下效率最高。
\end{proposition}
\admitted

\begin{example}[实验台版本]\label{ex:g7:digestion-and-nutrients:bench}
经典实验在玻璃管中进行\omterm{def:g4:journey-of-food:digestion}{消化}。两支试管
各盛有熟淀粉糊，置于体温\omterm{def:g6:exploring-our-environment:conditions}{条件}下：一支加入新鲜的
\omterm{def:g4:journey-of-food:tract}{唾液}，另一支加入煮沸（失效）的\omterm{def:g4:journey-of-food:tract}{唾液}。一小时后，
用淀粉检验发现2号管没有变化——仍是淀粉——
而1号管中的淀粉糊已变成糖，正是
即\cref{ex:g4:journey-of-food:sweet}中的那种甜味，
在没有口腔的情况下产生。单靠\omterm{def:g4:journey-of-food:digestion}{消化液}就完成了拆解——
而煮沸的对照表明，其有效成分
被热破坏，就像精密机器会被损坏一样。
\end{example}

\begin{remark}[为什么还要拆解？]\label{rem:g7:digestion-and-nutrients:why}
食物中的大分子进不去肠壁，
就像衣柜穿不过信箱一样。而人体
也不需要它们保持整体：它需要的是\emph{小块}——
小单位，它可以再组装成自身的物质
（见\cref{prop:g6:matter-from-food:transform}）。\omterm{def:g4:journey-of-food:digestion}{消化}既是为运输而拆解，
\emph{也}是为重组而拆解：
就像把家具拆成平板包装，每顿饭都在进行。
\end{remark}

\section{营养物质}

\begin{definition}[营养物质的类别]\label{def:g7:digestion-and-nutrients:families}
拆解产生的是\emph{营养物质}\index{营养物质}
——即可以直接进入\omterm{prop:g4:journey-of-food:absorb}{血液}的小分子成分：
\begin{itemize}
  \item \emph{糖类}（首先是\omterm{prop:g7:organs-at-work:bill}{葡萄糖}，
        即\cref{prop:g7:organs-at-work:bill}中的燃料），来自
        富含淀粉和甜味的食物；
  \item \emph{构建物质}的小单位，来自
        肉类、\omterm{prop:g3:sorting-living-things:groups}{鱼类}、蛋类、奶制品和豆类——这些是
        \cref{prop:g3:balanced-diet:jobs}所说
        “建造”工作的砖块；
  \item \emph{脂肪}的小单位——既是高密度的
        燃料，兼作建筑材料；
  \item 另外还有无需拆解的物质：
        \emph{维生素}、\emph{\omterm{prop:g4:what-plants-need:needs}{无机盐}}，以及水
        本身。
\end{itemize}
\cref{def:g3:balanced-diet:families}中的\omterm{def:g3:balanced-diet:families}{食物类别}
归根结底是这些营养物质类别的供应者——
“购物清单背后的购物清单”。
\end{definition}

\begin{example}[一顿午餐，逐项列出]\label{ex:g7:digestion-and-nutrients:lunch}
\cref{ex:g3:balanced-diet:lunch}中的均衡午餐，
按\omterm{def:g7:digestion-and-nutrients:families}{营养物质}重新解读：意面——淀粉，被拆解为
以升计的\omterm{prop:g7:organs-at-work:bill}{葡萄糖}；鱼——构建物质，切成其
小单位；豆子上的油——脂肪的小单位；豆子与苹果
——维生素、\omterm{prop:g4:what-plants-need:needs}{无机盐}、水，外加膳食纤维。注意，膳食纤维是
这是证明系统规律的例外：我们的任何\omterm{def:g4:journey-of-food:digestion}{消化液}都无法
切开它，因此它穿越不了任何壁，而是继续前行——
残渣的体积（见\cref{prop:g4:journey-of-food:absorb}），
承担着它自己有益的“清扫”工作。
\end{example}

\section{伟大的穿越}

\begin{proposition}[绒毛处的吸收]\label{prop:g7:digestion-and-nutrients:absorption}
前文承诺的穿越，
即\cref{prop:g4:journey-of-food:absorb}所述，发生在
\omterm{def:g4:journey-of-food:tract}{小肠}的内壁上，而这面墙正是为此而设计的：
有褶皱，并密布着数以百万计的
\emph{绒毛}\index{绒毛}——指状的小突起，每个
约一毫米高，各自包裹着自己的细小\omterm{def:g4:heart-and-blood:vessels}{血管}。
正是\omterm{prop:g7:digestion-and-nutrients:absorption}{绒毛}使得穿越面积达到
相当于铺满一个房间的规模（见\cref{rem:g4:journey-of-food:surface}）：
层层褶皱，再叠上指状\omterm{prop:g7:digestion-and-nutrients:absorption}{绒毛}。\omterm{def:g7:digestion-and-nutrients:families}{营养物质}穿过其薄壁
进入\omterm{prop:g4:journey-of-food:absorb}{血液}，\omterm{prop:g4:journey-of-food:absorb}{血液}首先把它们运送到
一个重要的分拣与储存器官——
\emph{肝脏}\index{肝脏}——然后才被分送到全身。
\end{proposition}
\admitted

\begin{omfigure}
\begin{tikzpicture}[scale=0.9]
  % intestine wall with villi
  \draw[thick] (0,0) -- (0,3.4);
  \draw[thick] (0,0) -- (8.2,0);
  \node[font=\small, rotate=90] at (-0.4,1.7) {小肠壁};
  % villi
  \foreach \x in {0.9,2.4,3.9,5.4,6.9} {
    \draw[thick, fill=omDef!8] (\x,0)
      .. controls ({\x-0.42},1.4) and ({\x-0.42},2.4) .. (\x,2.6)
      .. controls ({\x+0.42},2.4) and ({\x+0.42},1.4) .. (\x,0);
    % vessels inside
    \draw[omThm, thick] ({\x-0.14},0.15)
      .. controls ({\x-0.18},1.3) .. (\x,2.25)
      .. controls ({\x+0.18},1.3) .. ({\x+0.14},0.15);
  }
  \node[font=\footnotesize, align=center] at (4.1,-0.6)
    {每个绒毛包裹着自己的细小血管};
  % nutrients crossing
  \foreach \p in {(1.7,2.9),(3.2,3.0),(4.7,2.85),(6.2,2.95)} {
    \fill[omMeth] \p circle (0.07);
  }
  \node[font=\footnotesize, omMeth] at (6.9,3.25)
    {肠腔中的营养物质};
  \draw[->, omMeth, thick] (3.2,2.9) -- (3.65,2.2);
  \node[font=\footnotesize, align=left] at (10.0,1.7)
    {穿越过程：\\ 穿过绒毛的 \\ 薄壁，\\ 进入血液};
\end{tikzpicture}

{\small 肠道内壁的“地毯”：数以百万计的\omterm{prop:g7:digestion-and-nutrients:absorption}{绒毛}，
每根都是布满\omterm{def:g4:heart-and-blood:vessels}{血管}的薄壁“指状突起”。表面
再叠表面——这是人体最大的物质交换边界。}
\end{omfigure}

\begin{example}[再谈这一规格]\label{ex:g7:digestion-and-nutrients:spec}
大、薄、湿润、布满\omterm{def:g4:heart-and-blood:vessels}{血管}：\omterm{prop:g7:digestion-and-nutrients:absorption}{绒毛}恰好符合
\cref{prop:g7:how-animals-breathe:spec}所规定的规格——
只是交换的对象由\omterm{def:g7:what-is-respiration:respiration}{氧气}换成\omterm{def:g7:digestion-and-nutrients:families}{营养物质}。肺泡、
鳃丝、肠道\omterm{prop:g7:digestion-and-nutrients:absorption}{绒毛}：凡是身体需要大量穿越某一边界的地方，
就会建造同一种“机器”。认出这一规格一次，
便能处处识别它。
\end{example}

\begin{method}[追踪一种营养物质的全程]\label{met:g7:digestion-and-nutrients:trace}
对任意一口食物，追踪其中一种\omterm{def:g7:digestion-and-nutrients:families}{营养物质}：
\begin{enumerate}
  \item 说出该食物所属的类别及其主要成分；
  \item 逐站记录其分解的开始与结束位置，
        以及哪些\omterm{def:g4:journey-of-food:digestion}{消化液}参与了分解；
  \item 在\omterm{prop:g7:digestion-and-nutrients:absorption}{绒毛}处穿过肠壁进入\omterm{prop:g4:journey-of-food:absorb}{血液}，途经
        \omterm{prop:g7:digestion-and-nutrients:absorption}{肝脏}；
  \item 输送：作为燃料送往工作中的\omterm{def:g5:first-look-at-cells:cell}{细胞}，
        （\cref{prop:g7:organs-at-work:bill}），或
        作为原料送往建造部位
        （\cref{def:g6:matter-from-food:production}）。
\end{enumerate}
\end{method}

\begin{example}[以面包追踪营养]\label{ex:g7:digestion-and-nutrients:breadtrace}
面包：淀粉类，主要成分为淀粉。分解始于
口腔（\omterm{def:g4:journey-of-food:tract}{唾液}——细嚼后的回甘），
在胃酸中暂停，最终在\omterm{def:g4:journey-of-food:tract}{小肠}的\omterm{def:g4:journey-of-food:digestion}{消化液}中完成
分解为\omterm{prop:g7:organs-at-work:bill}{葡萄糖}。穿越过程：\omterm{prop:g7:digestion-and-nutrients:absorption}{绒毛}、\omterm{prop:g4:journey-of-food:absorb}{血液}、
\omterm{prop:g7:digestion-and-nutrients:absorption}{肝脏}——储存部分\omterm{prop:g7:organs-at-work:bill}{葡萄糖}以备后用，
这是它的众多功能之一。输送：大腿\omterm{def:g3:bones-and-muscles:muscle}{肌肉}在短跑途中，
一小时内即可将其燃烧。四年级那一口，
至此完全交代清楚。
\end{example}

\section{练习}

\begin{exercise}[$\star$]\label{exo:g7:digestion-and-nutrients:1}
\omterm{def:g4:journey-of-food:digestion}{消化液}能做到、而牙齿与胃的搅动
做不到的，是什么？
\end{exercise}

\begin{exercise}[$\star$]\label{exo:g7:digestion-and-nutrients:2}
描述双管\omterm{def:g4:journey-of-food:tract}{唾液}实验：装置、结果，
以及两条结论。
\end{exercise}

\begin{exercise}[$\star$]\label{exo:g7:digestion-and-nutrients:3}
列出各类\omterm{def:g7:digestion-and-nutrients:families}{营养物质}，并各举一种食物来源。
\end{exercise}

\begin{exercise}[$\star$]\label{exo:g7:digestion-and-nutrients:4}
哪些食物成分无需分解？哪一种
完全不能被\omterm{def:g4:journey-of-food:digestion}{消化液}分解？它有什么用处？
\end{exercise}

\begin{exercise}[$\star$]\label{exo:g7:digestion-and-nutrients:5}
什么是\omterm{prop:g7:digestion-and-nutrients:absorption}{绒毛}？每根\omterm{prop:g7:digestion-and-nutrients:absorption}{绒毛}包裹着什么？
\end{exercise}

\begin{exercise}[$\star$]\label{exo:g7:digestion-and-nutrients:6}
肠道的\omterm{prop:g4:journey-of-food:absorb}{血液}在被输送至全身之前先流向何处，
并说出该器官承担的一项功能。
\end{exercise}

\begin{exercise}[$\star\star$]\label{exo:g7:digestion-and-nutrients:7}
解释
\cref{rem:g7:digestion-and-nutrients:why}中“平板包装”的比喻，
即食物必须被分解的两个原因。
\end{exercise}

\begin{exercise}[$\star\star$]\label{exo:g7:digestion-and-nutrients:8}
煮沸\omterm{def:g4:journey-of-food:tract}{唾液}的试管有何意义？说出所用方法，
以及煮沸实验说明了\omterm{def:g4:journey-of-food:digestion}{消化液}中起作用的物质
具有什么性质。
\end{exercise}

\begin{exercise}[$\star\star$]\label{exo:g7:digestion-and-nutrients:9}
将\omterm{prop:g7:digestion-and-nutrients:absorption}{绒毛}与肺泡、鳃丝并列比较，
说明三者共有的规格，以及各自
所转运的物质。
\end{exercise}

\begin{exercise}[$\star\star$]\label{exo:g7:digestion-and-nutrients:10}
用\cref{met:g7:digestion-and-nutrients:trace}追踪一个煮鸡蛋（构建物质）
的营养，将其送往
愈合中擦伤处的修复部位（见\cref{exo:g6:cell-unit-of-life:10}）。
\end{exercise}

\begin{exercise}[$\star\star$]\label{exo:g7:digestion-and-nutrients:11}
极端节食连续数周跳过所有构建物质类食物。
结合\cref{def:g7:digestion-and-nutrients:families}与
\cref{def:g6:matter-from-food:production}，预测哪些物质会
首先短缺，以及在仍在生长的身体中，
什么最先受到影响。
\end{exercise}

\begin{exercise}[$\star\star\star$]\label{exo:g7:digestion-and-nutrients:12}
“\omterm{def:g4:journey-of-food:digestion}{消化}把一餐变成一份输送清单。”请
从头到尾为这句话辩护——\omterm{def:g4:journey-of-food:digestion}{消化液}、\omterm{def:g7:digestion-and-nutrients:families}{营养物质}、\omterm{prop:g7:digestion-and-nutrients:absorption}{绒毛}、
\omterm{prop:g7:digestion-and-nutrients:absorption}{肝脏}与输送各环节，用一段话予以论证。
\end{exercise}

\section{问题：玻璃肠道}

\begin{problem}[{周末问题——在实验台上搭建消化过程，
一管接一管}]\label{pb:g7:digestion-and-nutrients:1}
全班搭建一个“玻璃肠道”：一排体温\omterm{def:g6:exploring-our-environment:conditions}{条件}下的试管，
每支试管模拟\omterm{def:g4:journey-of-food:tract}{消化道}的一个站点。试管 M
（口腔）：淀粉糊加新鲜\omterm{def:g4:journey-of-food:tract}{唾液}。试管 M$'$：
淀粉糊加煮沸过的\omterm{def:g4:journey-of-food:tract}{唾液}。试管 S（胃）：蛋清块
加药房买来的胃液。试管
S$'$：同样的蛋清块放在清水中。四管均运行两小时，
随后进行检测：检测淀粉、检测糖类，
并观察蛋清块。

\medskip\noindent\textbf{第一部分——预测与
结果。}
\begin{enumerate}
  \item 预测四管两小时后的结果。
  \item M 管变甜，M$'$ 管不变；请给出结论，
        并说明 M$'$ 的确切作用。
  \item S 管中的蛋清块缩小、变软并变透明，
        S$'$ 管中的蛋清块保持不变。胃液做了什么，
        它产生的是哪一类
        \omterm{def:g7:digestion-and-nutrients:families}{营养物质}？
  \item 没有任何一管经过咀嚼或搅拌。
        这一排试管说明，\omterm{def:g4:journey-of-food:digestion}{消化}的
        主要工作发生在哪里？
\end{enumerate}

\medskip\noindent\textbf{第二部分——改进模型。}
\begin{enumerate}[resume]
  \item 有学生提议把 M 管放在冰箱\omterm{def:g6:exploring-our-environment:conditions}{温度}下运行，
        “为了整齐”。请据此提出异议，
        并把 \cref{prop:g6:decomposers-and-soil:speed} 的规律
        用于\omterm{def:g4:journey-of-food:digestion}{消化液}。
  \item 另一名学生提议用一支大试管同时放入两种\omterm{def:g4:journey-of-food:digestion}{消化液}，
        使其混合。
        \cref{prop:g7:digestion-and-nutrients:chemical} 中哪一条按站点区分的事实
        会被这种混合模糊？
  \item 玻璃肠道没有管壁，因此没有任何物质被吸收。
        完整的模型接下来需要哪种结构？
        它的材料必须具备哪两种性质？
  \item 在第五支试管中加入纤维（麦麸片），并提供全部\omterm{def:g4:journey-of-food:digestion}{消化液}。
        预测它两小时后的状态，
        以及它对模型的意义。
\end{enumerate}

\medskip\noindent\textbf{第三部分——从玻璃到肠道。}
\begin{enumerate}[resume]
  \item 在体内，M 管中的糖类接下来会遇到\omterm{prop:g7:digestion-and-nutrients:absorption}{绒毛}。
        请按四个站点描述它们的去向。
  \item 胃液在强酸中起作用；肠道中的\omterm{def:g4:journey-of-food:digestion}{消化液}
        则需要先把酸中和。
        这说明“每种\omterm{def:g4:journey-of-food:digestion}{消化液}都在其自身站点的\omterm{def:g6:exploring-our-environment:conditions}{条件}下起作用”
        这一说法意味着什么？以及按顺序吞下
        各个站点又说明了什么？
  \item 某制药公司出售\omterm{def:g4:journey-of-food:digestion}{消化液}胶囊，供自身站点
        功能不足的人服用。请结合试管
        实验说明，为什么这样的疗法
        究竟为何可行？
  \item 用一句话总结本次实验，
        句中须包含\emph{\omterm{def:g4:journey-of-food:digestion}{消化液}}、\emph{\omterm{def:g7:digestion-and-nutrients:families}{营养物质}}和
        \emph{对照管}。
\end{enumerate}
\end{problem}
